% Generated draft: requires mathematical and source review.
\documentclass[11pt]{article}
\usepackage[margin=25mm]{geometry}
\usepackage{fontspec}
\setmainfont{Latin Modern Roman}
\usepackage{amsmath,amssymb,mathtools,unicode-math}
\setmathfont{Latin Modern Math}
\usepackage{enumitem,needspace,xurl,hyperref}
\setlength{\parindent}{0pt}
\setlength{\parskip}{5pt}
\setlength{\emergencystretch}{3em}
\newcommand{\N}{\mathbb{N}}
\newcommand{\Z}{\mathbb{Z}}
\newcommand{\Q}{\mathbb{Q}}
\newcommand{\R}{\mathbb{R}}
\newcommand{\C}{\mathbb{C}}
\newcommand{\F}{\mathbb{F}}
\newcommand{\A}{\mathbb{A}}
\newcommand{\PP}{\mathbb{P}}
\newcommand{\E}{\mathbb{E}}
\newcommand{\eps}{\varepsilon}
\newcommand{\dd}{\,\mathrm{d}}
\begin{document}
\section*{UnsolvedMath ID: 700217}
\textbf{Source problem code:} catalogue entry
\subsection{Definitions and mathematical statement}
For a group \(G\), let \(G'=[G,G]\) be its commutator subgroup and let \(G''=[G',G']\). The variety \(\mathcal M\) of metabelian groups consists of all groups \(G\) satisfying \(G''=1\), with group homomorphisms as morphisms. A group \(P\in\mathcal M\) is projective in this variety if, for every surjective homomorphism \(q:H\twoheadrightarrow K\) between groups in \(\mathcal M\) and every homomorphism \(f:P\to K\), there exists a homomorphism \(\widetilde f:P\to H\) such that \(q\circ\widetilde f=f\).

For a set \(X\), let \(F(X)\) be the abstract free group on \(X\). The free metabelian group on \(X\) is \(F(X)/F(X)''\), the free object on \(X\) in \(\mathcal M\).

\textbf{Question (P. Hall; proposed by V. A. Artamonov).} Is every projective object of \(\mathcal M\) free in \(\mathcal M\)? Explicitly, is it true that
\[
\forall P\in\mathcal M,\qquad
P\text{ projective in }\mathcal M
\ \Longrightarrow\ 
\exists\text{ a set }X\text{ with }P\cong F(X)/F(X)''?
\]
No finiteness or rank restriction is supplied; the formulation allows arbitrary sets \(X\), including the empty set. Here the source's word ``free'' is interpreted as free in the variety of metabelian groups. The supplied wording does not explicitly distinguish this meaning from freeness as an abstract group, so that interpretation requires review.
\subsection{Short English statement}
Must every metabelian group that is projective within the variety of metabelian groups be a free metabelian group? The supplied source imposes no finiteness restriction. The meaning of ``free'' is interpreted relative to the variety and requires confirmation.
\subsection{Sources}
\begin{itemize}
\item Local UnsolvedMath registry, permanent ID \(700217\), exact source problem code \texttt{catalogue entry}, title \emph{Kourovka Notebook Problem 4.6}. The supplied registry labels the problem open; this label has not been independently verified.
\item \emph{The Kourovka Notebook (Unsolved Problems in Group Theory)}, original problem 4.6, fourth issue (1973); supplied edition reference: twenty-first issue edition (2026), arXiv:1401.0300v48, \url{https://arxiv.org/abs/1401.0300v48}. This reference is supplied evidence and has not been independently consulted.
\item Attribution in the supplied statement: P. Hall. The proposer field and background identify V. A. Artamonov, with proposed year 1973. The supplied collection description names E. I. Khukhro and V. D. Mazurov as editors.
\end{itemize}
\end{document}
